\section{Consequences for the endpoints}
\label{sec:endpoints}

Theorem~\ref{thm:main} makes every maximal program finite.  The following
consequences use log abundance and bounded complements to describe its
endpoint more precisely.  The semiample-endpoint statement holds over
every algebraically closed field of characteristic zero; the uniform
local statements are over $\mathbb C$.

\begin{corollary}[Semiample endpoints]\label{cor:semiample-endpoints}
Let $k$ be an algebraically closed field of characteristic zero, and let
$(X,B)$ be a projective lc fourfold pair over $k$ with rational boundary.
Suppose that $K_X+B$ is pseudo-effective. Then every maximal program of
Definition~\ref{def:program} ends at a projective lc pair $(Y,B_Y)$ for
which some positive Cartier multiple of $K_Y+B_Y$ is generated by global
sections.
\end{corollary}
\begin{proof}
Fix a maximal program $X_0\dashrightarrow\cdots\dashrightarrow X_r$, finite
by Theorem~\ref{thm:main}, and put $D_i=K_{X_i}+B_i$.
The section-space comparison in
\cite[proof of Corollary~12.1(i)]{OpenAILogAbundance2026}
identifies the sections of the pseudo-effective $D_0$ with those of a
semiample birational adjoint.  Hence $D_0$ has a nonzero section in some
positive Cartier degree.  Choose a common positive multiple $N$ of that
degree and the Cartier indices of $D_0,\ldots,D_r$.  A suitable power of
the section is represented by a rational function $s$ satisfying
$\operatorname{div}_{X_0}(s)+ND_0\geq0$.  At a birational step, push the
effective divisor $\operatorname{div}_{X_i}(s)+ND_i$ through
$f_i$ and take its strict transform through $f_i^+$. The boundary rule and
compatible canonical divisors, together with smallness of $f_i^+$, identify
the result with $\operatorname{div}_{X_{i+1}}(s)+ND_{i+1}\geq0$. Thus the
section passes through every divisorial, flipping or mixed step.

A Mori fibre endpoint has a positive-dimensional fibre meeting the complement
of this effective Cartier zero divisor; choose an integral curve $C$ in that
fibre through such a point. Its intersection with the zero divisor is
nonnegative, but equals $ND_r\cdot C<0$ by relative ampleness of $-D_r$, a
contradiction. The endpoint is therefore nef, and log abundance
\cite[Theorem~1.1]{OpenAILogAbundance2026} makes its rational adjoint semiample.
\end{proof}

For $\epsilon>0$, a pair is \emph{$\epsilon$-lc} if every normalized
divisorial valuation has log discrepancy at least $\epsilon$.
A projective contraction $h:Y\to S$ is \emph{of Fano type} if some
effective rational boundary $\Gamma$ on $Y$ makes $(Y,\Gamma)$ klt and
$-(K_Y+\Gamma)$ ample over $S$.  The next consequence uses this condition at a Mori
endpoint.

\begin{corollary}[Uniform local complements at Mori endpoints]
\label{cor:uniform-mori-complements}
Fix a real number $\epsilon>0$ and a finite set $I\subset[0,1]\cap\mathbb Q$.
There is a positive integer $M=M(4,\epsilon,I)$ with the following
property. Let $(X,B)$ be a projective $\epsilon$-lc fourfold pair over
$\mathbb C$ with coefficients in $I$ and with $K_X+B$ not pseudo-effective.
Every maximal program of Definition~\ref{def:program} ends at a Mori
contraction $h:(Y,B_Y)\to S$. For every closed $s\in S$, there is an
effective $\mathbb Q$-divisor $B_Y^+\geq B_Y$ such that, on the inverse
image of some neighbourhood of $s$, the pair $(Y,B_Y^+)$ is klt and
$M(K_Y+B_Y^+)$ is Cartier and linearly trivial.
\end{corollary}
\begin{proof}
Theorem~\ref{thm:main} makes the program finite.
Lemma~\ref{lem:graph} shows that log discrepancies do not decrease,
including at mixed steps, while the boundary rule and smallness of the
positive maps retain or delete the original coefficients.
On a common smooth projective model $p:W\to X$, $q:W\to Y$, the graph
comparisons give
$p^*(K_X+B)=q^*(K_Y+B_Y)+E$ with $E\geq0$.
If the final adjoint were nef, this identity and birational invariance
of pseudo-effectivity for rational Cartier divisors would make
$K_X+B$ pseudo-effective. The endpoint is therefore Mori.
It remains $\epsilon$-lc, hence klt, and $B_Y$ itself witnesses that
$Y$ is of Fano type over $S$, since $-(K_Y+B_Y)$ is $h$-ample.
The varieties are projective, so
\cite[Corollary~1.2]{OpenAIComplements2026} applies and gives the claimed
index, independent of the program and its endpoint.
\end{proof}

Under the hypotheses of Corollary~\ref{cor:uniform-mori-complements},
when $\dim S>0$, \cite[Theorem~1]{OpenAICartier2026} also gives a uniform
$\tau=\tau(4,\epsilon)>0$ such that, for each closed $s\in S$, some
neighbourhood $U\ni s$ carries a nonzero effective Cartier divisor $T$
through $s$ with
$(h^{-1}(U),B_Y|_{h^{-1}(U)}+\tau h^*T)$ log canonical.

